\section{Load both inputs before an in-place store}\label{pcm-aliases}
When the output equals the first input, each stored result replaces a sample that
was just read from that input. The same rule applies when output equals the second
input. Both values must be retained before the first write to their shared sample
position. Loading only A and then writing a provisional output could destroy B
before its load when the destination is B.

A complete vector group follows the same reasoning at a larger width. Load both
sixteen-byte input groups into registers before storing their sixteen-byte output.
At an exact alias, the next group begins after these bytes, so the store does not
alter a future group's inputs. Shifted overlap would break that argument; it is
excluded rather than silently treated as a supported copy direction.

\begin{figure}[htbp]
\centering
\begin{tikzpicture}[alt={Two complete input groups, A and B, are loaded into distinct registers. Both arrows reach an addition stage before a separate arrow reaches the output store. The output can be A, B or a separate slice. An input load must not be delayed until after a store that can overwrite it.},node distance=13mm and 19mm,>=Stealth]
\node[draw,align=center] (a) {Load A group\\16 bytes};
\node[draw,align=center,right=of a] (b) {Load B group\\16 bytes};
\node[draw,align=center,below=of b] (sum) {Add and count\\using loaded values};
\node[draw,align=center,left=of sum] (out) {Store output\\16 bytes};
\draw[->] (a)--(sum);
\draw[->] (b)--(sum);
\draw[->] (sum)--(out);
\end{tikzpicture}
\caption{An exact in-place alias changes the destination's identity, while the dependency on both input values remains.}
\end{figure}

Logical bounds still apply to every access. Eight samples make sixteen bytes, so
both vector loads and the store fit when at least eight samples remain. A shorter
tail uses byte-based scalar operations. Padding in the 192-byte backing arrays is
not an extension of a passed slice.
